\documentclass[11pt]{article}
\usepackage{amsmath,amssymb}
\usepackage[margin=1in]{geometry}
\usepackage{siunitx}

\title{On the normalization of Eqs.~(7)--(9) of Hsu \& Lai,\\
\emph{Angle-Resolved Spectroscopy of Parametric Fluorescence}\\
\large (arXiv:1302.6226) --- angle conventions and the genuine misprints}
\author{}
\date{}

\newcommand{\sinc}{\operatorname{sinc}}
\newcommand{\dd}{\mathrm{d}}

\begin{document}
\maketitle

\noindent
The paper's master formula for the angle-resolved Type-I SPDC photon flux
(used to compute Fig.~6) is Eq.~(9). An earlier version of this note claimed
that the printed Eq.~(9) was $2\pi n_s^2 \approx 17\times$ too low, attributing
the deficit to a dropped coefficient in the angular Jacobian. That claim was
\textbf{wrong}: it rested on reading $\theta_s$ as the \emph{internal}
(in-crystal) angle. The paper's Fig.~2 defines the angle conventions
explicitly --- ``The angles $\theta_s'$ and $\theta_s$ are, respectively,
internal and external angles formed by the signal and pump wave vectors'' ---
and with $\theta_s$ the \emph{external} (laboratory) angle the printed
Eq.~(9) is exactly correct as a flux density per unit azimuth. Its only
blemish is a dropped $\dd\phi$ symbol. The genuine misprints are confined to
the main-text Eq.~(7): the prefactor constant ($2\pi^4$ vs the appendix's
correct $8\pi^4$) and the sign of the Gaussian exponent.

\section*{The relevant printed equations}

\paragraph{Eq.~(7), main text --- flux per unit frequency:}
\begin{equation}
N_s(\omega_s,\kappa_s) =
\frac{\hbar d_{\mathrm{eff}}^2\,\omega_s\omega_i\omega_p L^2 N_p}
{\mathbf{2\pi^4}\, c^3\varepsilon_0 n_s n_i n_p}\,
\dd\omega_s\,\dd^2\kappa_s
\int \dd^2\xi\,
e^{+\frac12\xi^2}\,
\sinc^2\!\left(\tfrac12 L\,\Delta k_z\right).
\tag{7}
\end{equation}
(Both bold-faced items are misprints; see below. The exponent must be
$-\tfrac12\xi^2$, as the appendix has it.)

\paragraph{Appendix intermediate} (after relabeling $\xi\to\xi_i$ and the
saddle-point integral over the \emph{relative} azimuth between
$\kappa_s$ and $\kappa_i$, which supplies the $\sqrt{2\pi}$):
\begin{equation}
N_s =
\underbrace{\frac{\hbar d_{\mathrm{eff}}^2\,\omega_s\omega_i\omega_p L^2 N_p\,\sqrt{2\pi}}
{\mathbf{8\pi^4}\, c^3\varepsilon_0 n_s n_i n_p}}_{\textstyle A_0}\,
\dd\omega_s\,\underbrace{\kappa_s\,\dd\kappa_s\,\dd\phi}_{=\,\dd^2\kappa_s}\,
\int \dd\xi_i\, e^{-\frac12(\xi_s-\xi_i)^2}\,
\sinc^2\!\left(\tfrac12 L\,\Delta k_z\right).
\label{eq:inter}
\end{equation}
Note the signal azimuth $\dd\phi$ is still present here.

\paragraph{Eq.~(9), as printed} (target: variables $\lambda_s,\theta_s$):
\begin{equation}
N_s(\lambda_s,\theta_s) =
2\pi\sqrt{2\pi}\,
\frac{\hbar d_{\mathrm{eff}}^2\,\omega_p L^2 N_p}
{\varepsilon_0 n_s n_i n_p\,\lambda_s^5\lambda_i}\,
\sin(2\theta_s)\,\dd\lambda_s\,\dd\theta_s
\int \dd\xi_i\, e^{-\frac12(\xi_s-\xi_i)^2}\,
\sinc^2\!\left(\tfrac12 L\,\Delta k_z\right).
\tag{9}
\end{equation}

\section*{The angle convention: $\theta_s$ is the laboratory angle}

The paper primes all in-medium quantities: Eq.~(3) and Eq.~(6) use the
internal angle $\theta_s'$, and Eq.~(6) and the appendix define the in-medium
wavenumbers as $k_s' = n_s\omega_s/c$, etc. Accordingly, the appendix's
substitution ``$\kappa_s = k_s\sin\theta_s$'' (unprimed $k_s$, unprimed
$\theta_s$) means the \emph{vacuum} wavenumber $k_s = k_0 = 2\pi/\lambda_s$
with the \emph{external} angle. This is also the physically correct statement
of the conserved transverse momentum: $\kappa$ is continuous across the flat
exit face (Snell), so
\begin{equation}
\kappa_s \;=\; k_0\sin\theta_s^{\text{ext}}
        \;=\; k_s'\sin\theta_s'^{\,\text{int}},
\qquad k_0 = \frac{2\pi}{\lambda_s},\quad k_s' = \frac{2\pi n_s}{\lambda_s}.
\end{equation}
Both parametrizations describe the same $\kappa_s$; the external one carries
no $n_s$.

\section*{Deriving Eq.~(9) from the appendix (external angle)}

Carry \eqref{eq:inter} to $(\lambda_s,\theta_s)$ via
\begin{align}
\textbf{(S1)}\quad
\dd\omega_s &= \frac{2\pi c}{\lambda_s^2}\,\dd\lambda_s
&&\left(\omega_s=\tfrac{2\pi c}{\lambda_s}\right),\\[4pt]
\textbf{(S2)}\quad
\kappa_s\,\dd\kappa_s
&= \tfrac12 k_0^2\,\sin(2\theta_s)\,\dd\theta_s
= \frac{2\pi^2}{\lambda_s^2}\,\sin(2\theta_s)\,\dd\theta_s
&&\left(\kappa_s=k_0\sin\theta_s\right),
\end{align}
with $\omega_s\omega_i = 4\pi^2 c^2/(\lambda_s\lambda_i)$ inside $A_0$ and the
signal azimuth $\dd\phi$ kept as a differential. The numeric coefficient is
\begin{equation}
\frac{1}{8\pi^4}\cdot 4\pi^2 \cdot 2\pi \cdot 2\pi^2 = 2\pi ,
\end{equation}
so
\begin{equation}
\boxed{\;
N_s(\lambda_s,\theta_s,\phi) =
2\pi\sqrt{2\pi}\,
\frac{\hbar d_{\mathrm{eff}}^2\,\omega_p L^2 N_p}
{\varepsilon_0\, n_s n_i n_p\,\lambda_s^5\lambda_i}\,
\sin(2\theta_s)\,\dd\lambda_s\,\dd\theta_s\,\dd\phi
\int \dd\xi_i\, e^{-\frac12(\xi_s-\xi_i)^2}\,
\sinc^2\!\left(\tfrac12 L\,\Delta k_z\right)
\;}
\label{eq:corrected}
\end{equation}
which is the printed Eq.~(9) \emph{verbatim}, apart from the $\dd\phi$. The
printed equation is therefore a flux density per unit azimuth; the full ring
carries an extra factor $\int_0^{2\pi}\dd\phi = 2\pi$. This is not a
numerical error in the paper: Sec.~III.D states explicitly that the flux is
integrated over ``$\phi = 0 \to 2\pi$'', so the $2\pi$ enters at the
numerical-integration stage. The only defect of Eq.~(9) is the missing
$\dd\phi$ symbol.

Here $\xi_s=\kappa_s W$ and
$\Delta k_z = \sqrt{k_s'^2-\kappa_s^2} + \sqrt{k_i'^2-(\xi_i/W)^2} - k_p'$
(appendix form; the in-medium wavenumbers stay primed). Integrating
\eqref{eq:corrected} over the azimuth reproduces
$\eta=2N_s/N_p=\num{2.72e-10}$ and $2N_s=\SI{4.44e7}{s^{-1}}$
($\theta_m=\SI{29.12}{\degree}$, $P_p=\SI{80}{mW}$), matching the paper's
Table~I. Implemented in \texttt{spdc\_physics.py} (\texttt{Ns} with
\texttt{external=True}; the equivalent internal-angle parametrization
$\kappa_s = k_s'\sin\theta_s'$, whose Jacobian carries $\tfrac12 k_s'^2 =
\tfrac12 n_s^2 k_0^2$, is available as \texttt{external=False} and yields the
identical total flux).

\paragraph{Where the old ``$2\pi n_s^2$'' claim came from.}
Reading ``$\kappa_s = k_s\sin\theta_s$'' with the \emph{in-medium}
$k_s = 2\pi n_s/\lambda_s$ (and hence $\theta_s$ internal) makes the Jacobian
(S2) acquire a factor $n_s^2$, and folding the azimuthal $2\pi$ into the
prefactor as well makes the printed Eq.~(9) appear short by
$2\pi n_s^2 = 2\pi(1.660)^2 \approx 17.3$. Under Fig.~2's convention neither
factor is missing: the $n_s^2$ never belonged there, and the $2\pi$ is the
azimuth integral the authors perform explicitly.

\section*{The genuine misprints: Eq.~(7)}

\paragraph{The factor of 4 ($2\pi^4$ vs $8\pi^4$).}
Eq.~(7) and \eqref{eq:inter} are the same equation up to the relabeling
$\int\dd^2\xi \to \int\dd^2\xi_i$ with $\xi=(\kappa_s+\kappa_i)W$,
$\xi_i=\kappa_i W$. At fixed $\kappa_s$ this is a pure \emph{shift},
\begin{equation}
\int \dd^2\xi\; f(\xi) \;=\; \int \dd^2\xi_i\; f(\xi_s+\xi_i),
\qquad \text{Jacobian}=1,
\end{equation}
which cannot alter the prefactor, so one of the two printed constants is a
misprint. The appendix's $8\pi^4$ is the correct one, on three independent
grounds: (i)~it reproduces the paper's simulated $\eta$
(Table~I: $\num{2.6e-10}$; our integral: $\num{2.72e-10}$);
(ii)~it agrees with the closed-form Eq.~(8) --- taken from the independent
Ref.~[18] (Koch et al., IEEE J.\ Quantum Electron.\ \textbf{31}, 769 (1995))
--- which evaluated directly gives $\eta = \num{2.63e-10}$;
(iii)~it matches the measured $\eta \approx$ \numrange{2.7e-10}{2.8e-10}.
With $2\pi^4$ everything would sit $4\times$ above all three.

\paragraph{The exponent sign.}
Eq.~(7) prints $e^{+\frac12\xi^2}$, which diverges; the appendix's
$e^{-\frac12|\xi_s+\xi_i|^2}$ is correct.

\paragraph{Summary.}
\begin{center}
\begin{tabular}{lll}
\hline
Item & Location & Status\\
\hline
$e^{+\frac12\xi^2}$ & Eq.~(7) exponent & misprint; must be $-$ (appendix has it right)\\
$2\pi^4$ & Eq.~(7) prefactor & misprint; correct value $8\pi^4$ (appendix, Eq.~(8), experiment)\\
missing $\dd\phi$ & Eq.~(9) & notation only; Eq.~(9) is a per-unit-azimuth density\\
``dropped $2\pi n_s^2$'' & Eq.~(9) & \textbf{retracted}: $\theta_s$ is the external angle (Fig.~2),\\
& & so the printed Jacobian coefficient is exact\\
\hline
\end{tabular}
\end{center}

\end{document}
